\chapterhead{46}{ORR-7}{Integrated Risk Management}{4}{0}

\lohead{46}{a}{Describe the role of risk governance, risk appetite, and risk
culture in the context of an enterprise risk management (ERM) framework.}

\orsub{Introduction to ERM}

\begin{ordefbox}{Enterprise Risk Management}
The ERM framework includes the tools and methods that an organisation uses to
manage the different types of risks that can impact the achievement of business
goals and objectives. This framework involves four key stages.
\end{ordefbox}

\orsubb{Risk management cycle (4 stages)}

\begin{center}
\begin{tikzpicture}[font=\fontsize{8.2}{10.2}\selectfont,
  st/.style={text width=3.2cm,align=center,rounded corners=3pt,inner sep=6pt,
             text=white,font=\sffamily\bfseries\fontsize{8.4}{10.2}\selectfont},
  sb/.style={text width=3.2cm,align=center,rounded corners=3pt,inner sep=6pt,
             fill=paleblue,draw=palenavy,line width=0.6pt,text=navy}]
\foreach \i/\x/\c/\t/\d in {%
  1/-6.0/navy/{1) Identification}/{Recognising potential risks},
  2/-2.0/teal/{2) Assessment}/{Evaluating the likelihood and impact of risks},
  3/2.0/forest/{3) Mitigation}/{Reducing the impact of identified risks},
  4/6.0/olive/{4) Reporting and monitoring}/{Tracking risks and the effectiveness of controls}}
{
  \node[st,fill=\c,anchor=north] (t\i) at (\x,0) {\t};
  \node[sb,anchor=north] (b\i) at (t\i.south) {\d};
}
\foreach \a/\b in {1/2,2/3,3/4}
  \draw[-{Stealth[length=4pt]},navy!35,line width=1.2pt]
    ($(t\a.east)+(0.04,0)$) -- ($(t\b.west)+(-0.04,0)$);
\end{tikzpicture}
\end{center}
\orfigcap{The four stages of the risk management cycle.}

\orsub{Structure of ERM}

\orsubb{Risk governance}

Sets roles and responsibilities for managing risk.

\begin{itemize}
\item \textbf{First line:} business line staff --- identify, measure, mitigate,
and report risks.
\item \textbf{Second line:} oversees risk management --- monitor risk types,
establish risk tools, and oversee processes.
\item \textbf{Third line:} independent reviewers (internal audit) --- evaluate
risk management effectiveness.
\end{itemize}

\orsubb{Risk culture}

Behaviours and values around managing risk. A strong risk culture leads to better
ERM effectiveness.

\orsubb{Risk appetite}

Defines the acceptable level of risk for achieving objectives. Banks use risk
appetite to set limits for credit risk (lending), market risk (trading), and
liquidity risk (cash flow).

% ---------------------------------------------------------------------
\lohead{46}{b}{Summarize the role of Basel regulatory capital and the process of
determining internal economic capital.}

\orsub{Capital elements of the ERM framework}

\orsubb{1) Regulatory capital}

Minimum capital levels mandated by the Basel Committee on Banking Supervision
(BCBS) to cover credit, market, operational risks, and liquidity (Basel I, II,
and III).

\noindent\begin{minipage}{\textwidth}
\renewcommand{\arraystretch}{1.35}
\begin{tabularx}{\textwidth}{@{}L{0.14\textwidth}Y@{}}
\rowcolor{navy} \thd{Pillar} & \thd{What it does}\\
{\tabfont \textbf{Pillar 1}} & {\tabfont Sets minimum capital requirements for market, credit, operational risks, and liquidity.}\\
\rowcolor{paleblue}
{\tabfont \textbf{Pillar 2}} & {\tabfont Allows supervisors to adjust capital requirements for individual institutions.}\\
{\tabfont \textbf{Pillar 3}} & {\tabfont Requires mandatory disclosure of risk information and financial health.}\\
\end{tabularx}
\end{minipage}
\orfigcap{The three pillars, as they apply to the capital framework.}

\orsubb{2) Economic capital}

The capital needed by an individual institution to cover unexpected losses based
on its specific risk profile, exceeding regulatory requirements. Credit rating
plays a role here.

\orsubb{3) Risk-Adjusted Return on Capital (RAROC)}

A metric used to assess profitability relative to the risks undertaken. It
considers expected after-tax net income, adjusted for potential losses, divided by
economic capital. Easier to apply for credit risks than market risks.

\orsubb{4) Capital assessment risk aggregation}

Strategies to determine total capital needs considering interactions between
risks. \textbf{Risk aggregation} combines capital needs for different risk
classes, considering that they won't all hit at once.

% ---------------------------------------------------------------------
\lohead{46}{c}{Describe elements of a stress-testing framework for financial
institutions and explain best practices for stress testing.}

Stress testing simulates how an institution would perform under extreme market
conditions. It is a crucial tool for assessing financial stability and informing
capital requirements.

Before the 2007--2009 crisis, stress testing was quantitative-focused. Post-crisis,
it now combines both qualitative and quantitative elements.

\orsub{Taxonomy of stress testing}

\begin{enumerate}
\item[\textcolor{rust}{\sffamily\bfseries 1.}] \textbf{Analytical method:}
qualitative (scenario analysis) to quantitative (parameter impact modelling).
\item[\textcolor{rust}{\sffamily\bfseries 2.}] \textbf{Risk type:} measurable
(probabilities assigned) to immeasurable (``unknown unknowns'').
\end{enumerate}

\begin{center}
\begin{tikzpicture}[font=\fontsize{7.8}{9.6}\selectfont,
  pt/.style={circle,fill=navy,inner sep=2pt},
  capt/.style={font=\sffamily\bfseries\fontsize{7.6}{9.4}\selectfont,text=white,
              align=center}]

\fill[palenavy] (-6.6,1.35) -- (-5.9,1.85) -- (-5.9,1.55) -- (5.9,1.55)
  -- (5.9,1.85) -- (6.6,1.35) -- (5.9,0.85) -- (5.9,1.15) -- (-5.9,1.15)
  -- (-5.9,0.85) -- cycle;
\node[capt,text=navy] at (0,1.35) {Dimension 1: Quantitative \ $\longleftrightarrow$ \ Qualitative approach};

\draw[midgray,line width=0.7pt] (-6.0,0) -- (6.0,0);
\node[pt] (p1) at (-4.2,0) {};
\node[pt] (p2) at ( 2.0,0) {};
\node[pt] (p3) at ( 4.6,0) {};
\node[text=navy,anchor=north,font=\sffamily\fontsize{7.4}{9}\selectfont] at (-4.2,-0.2) {Parameter stress testing};
\node[text=navy,anchor=north,font=\sffamily\fontsize{7.4}{9}\selectfont] at ( 2.0,-0.2) {Macro stress testing};
\node[text=navy,anchor=north,font=\sffamily\fontsize{7.4}{9}\selectfont] at ( 4.9,-0.75) {Reverse stress testing};

\fill[navy] (-6.6,-1.75) -- (-5.9,-1.25) -- (-5.9,-1.55) -- (5.9,-1.55)
  -- (5.9,-1.25) -- (6.6,-1.75) -- (5.9,-2.25) -- (5.9,-1.95) -- (-5.9,-1.95)
  -- (-5.9,-2.25) -- cycle;
\node[capt] at (0,-1.75) {Dimension 2: Measurable \ $\longleftrightarrow$ \ Immeasurable risk};
\end{tikzpicture}
\end{center}
\orfigcap{The stress-testing taxonomy, across two dimensions.}

\orsub{Stress testing approaches}

\noindent\begin{minipage}{\textwidth}
\renewcommand{\arraystretch}{1.35}
\begin{tabularx}{\textwidth}{@{}L{0.20\textwidth}L{0.24\textwidth}Y@{}}
\rowcolor{navy} \thd{Approach} & \thd{Where it sits} & \thd{What it does}\\
{\tabfont \textbf{a) Parameter testing}} & {\tabfont Quantitative, measurable} &
{\tabfont Changes model parameters to assess stress impacts.}\\
\rowcolor{paleblue}
{\tabfont \textbf{b) Macroeconomic testing}} & {\tabfont Quantitative and qualitative; measurable and immeasurable} &
{\tabfont Evaluates solvency using regulator-provided scenarios, stressing all risk types.}\\
{\tabfont \textbf{c) Reverse testing}} & {\tabfont Qualitative, immeasurable} &
{\tabfont Analyses ``unknown unknowns'' through scenario analysis, focusing on operational resilience and resolution planning.}\\
\end{tabularx}
\end{minipage}
\orfigcap{Three approaches, placed on the two taxonomy dimensions.}

% ---------------------------------------------------------------------
\lohead{46}{d}{Explain challenges and considerations when developing and
implementing models used in stress testing operational risk.}

\orsub{Stress testing operational risk}

Operational risk stress testing has evolved significantly, especially after the
financial crises and the COVID-19 pandemic. Stress testing has evolved from being
primarily quantitative to relying more on expert judgment.

The US CCAR sets the benchmark for operational risk stress testing expectations.

\begin{enumerate}
\item[\textcolor{rust}{\sffamily\bfseries 1.}] \textbf{Expected nonlegal loss
forecast module:} quantitative model plus expert refinement to estimate loss
forecasts for each risk type under different economic scenarios.
\item[\textcolor{rust}{\sffamily\bfseries 2.}] \textbf{Legal loss module:}
forecasts losses for litigation cases, considering the delay between
macroeconomic events and legal losses.
\item[\textcolor{rust}{\sffamily\bfseries 3.}] \textbf{Idiosyncratic scenario
add-on module:} captures unique risk exposures for each bank based on specific
storylines and vulnerabilities.
\end{enumerate}

\begin{ornotebox}{Remember}
Recognise the connection between operational risks and macroeconomic conditions,
despite some debate.
\end{ornotebox}

\orsub{Modeling approaches for operational risk}

\noindent\begin{minipage}{\textwidth}
\renewcommand{\arraystretch}{1.35}
\begin{tabularx}{\textwidth}{@{}L{0.20\textwidth}Y@{}}
\rowcolor{navy} \thd{Approach} & \thd{How it works}\\
{\tabfont \textbf{1) Preferred}} &
{\tabfont Model frequency and severity separately using regression analysis to
capture macroeconomic dependencies.}\\
\rowcolor{paleblue}
{\tabfont \textbf{2) Modified LDA}} &
{\tabfont Use the Loss Distribution Approach with Monte Carlo simulations, which
assumes unchanging risk exposures. A modified LDA can incorporate macroeconomic
factors for frequency.}\\
{\tabfont \textbf{3) Conditional LDA}} &
{\tabfont A balance between regression and LDA. Models frequency with regression
and uses expert judgment for severity. The challenge is choosing a realistic
severity percentile.}\\
\end{tabularx}
\end{minipage}
\orfigcap{Three modelling approaches for stress testing operational risk.}

\orsubb{Model refinement}

Because severity cannot be found reasonably by modelling alone, experts and risk
owners should review and potentially adjust the model process, inputs, outputs,
historical data, chosen approaches, and loss estimates.
